\clearpage
\chapter*{2024 Febrero}
\addcontentsline{toc}{chapter}{2024 Febrero}

\begin{enunciado}
    \hypertarget{2024 B 1}{}
    \textbf{1. (3,5 puntos)} Se considera el espacio topológico \((\mathbb{R}, T)\), donde
    \[
        T = \{\emptyset\} \cup \{M \subset \mathbb{R} \text{ tal que } \mathbb{Q} - M \text{ es un conjunto finito}\}.
    \]
    \begin{enumerate}[label=\alph*)]
        \item Hallar el interior, la adherencia y la frontera del conjunto \(\mathbb{R} - \mathbb{Q}\).
        \item Estudiar si \((\mathbb{R}, T)\) verifica el primer axioma de numerabilidad.
    \end{enumerate}
\end{enunciado}

\begin{enunciado}
    \hypertarget{2024 B 2}{}
    \textbf{2. (3,5 puntos)} Sea \((Y, S)\) un espacio de Hausdorff, \(f : (X, T) \to (Y, S)\) una aplicación y \(G = \{(x, f(x)) \mid x \in X\}\) el grafo de \(f\) en \(X \times Y\).
    \begin{enumerate}[label=\alph*)]
        \item Supongamos que \(f\) es continua. Estudiar si \(G\) es cerrado en el espacio topológico producto.
        \item Supongamos que \(G\) es cerrado en el espacio topológico producto. Estudiar si \(f\) es continua.
    \end{enumerate}
\end{enunciado}

\begin{enunciado}
    \hypertarget{2024 B 3}{}
    \textbf{3. (3 puntos)} Sea \((X, T)\) un espacio topológico y sea \(p \in X\) un punto.
    \begin{enumerate}[label=\alph*)]
        \item Definir la componente conexa de \(p\) en \((X, T)\).
        \item Estudiar si las componentes conexas de \((X, T)\) son conjuntos abiertos.
        \item Estudiar si las componentes conexas de \((X, T)\) son conjuntos cerrados.
    \end{enumerate}
\end{enunciado}